% =====================================================================
% Section 3: Problem Formulation
% =====================================================================
%
% Status: Draft Section 3 Problem Formulation v2 (2026-05-17)
% Revision note (v2 vs v1):
%   - Section title renamed from "Task Formulation" to
%     "Problem Formulation" per advisor feedback (align with §4 method
%     formulation language).
%   - Label renamed: sec:task -> sec:problem.
%   - Body unchanged; only the framing word changed.
%
% Purpose:
%   - Define the retrieval-first interface for EvidenceFlow
%   - Introduce ranked evidence list R_q and reader prefix
%   - Define document-grounded fact units and provenance
%   - Avoid Method details such as G_q, PPR, expansion, and noisy-OR
%
% =====================================================================

\section{Problem Formulation}
\label{sec:problem}

Given a natural language question $q$ and a corpus
$\mathcal{D}=\{d_i\}_{i=1}^{N}$, multi-hop question answering
requires evidence from multiple passages before producing an answer.
The retriever outputs an ordered evidence list
\[
\begin{aligned}
R_q &= (d_{(1)},d_{(2)},\ldots), \\
R_q^{(K)} &= (d_{(1)},\ldots,d_{(K)}), \\
\hat{a}_q &= \mathrm{LLM}(\mathrm{prompt}(q,R_q^{(K)})),
\end{aligned}
\]
and the reader consumes the prefix $R_q^{(K)}$. Retrieval quality
is evaluated by whether prefixes of $R_q$ contain the gold
supporting passages, while QA quality compares $\hat{a}_q$ with
the reference answer.

Pointwise relevance ranks each $d \in \mathcal{D}$ independently,
but multi-hop QA needs the prefix $R_q^{(K)}$ to be considered
\emph{as a set}: it should be both \emph{connected} across
documents, preserving bridge relations that span passages, and
\emph{covering}, in the sense that the selected $K$ passages
collectively address the distinct entities and relations $q$ asks
about. We therefore frame the retrieval problem as an evidence-set
objective alongside the ranked list,
\begin{equation}
\label{eq:evidence-set}
E_q^{\star}
=
\operatorname*{arg\,max}_{E \subseteq \mathcal{D},\,|E|=K}
\;\mathrm{ConnCov}(E\,;\,q),
\end{equation}
where $\mathrm{ConnCov}$ is a connected-and-covering criterion
that this work instantiates with the noisy-OR coverage score of
\S\ref{sec:method:enhancement}. The retrieval interface stays a
ranked list, but the prefix $R_q^{(K)}$ is judged by how close it
comes to a connected and covering set under
Eq.~\ref{eq:evidence-set}.

During indexing, each passage is associated with
document-grounded OpenIE facts. A fact unit is $f=(h,r,t)$, with
argument phrases $h, t$ and relation phrase $r$. Each fact
retains its provenance $\mathrm{src}(f)\in\mathcal{D}$, and
\[
\mathcal{F}_d=\{f\in\mathcal{F}: \mathrm{src}(f)=d\}
\]
denotes the facts extracted from passage $d$. The goal of
\methodname{} is to use these provenance-linked facts to rank
passages so that prefixes of $R_q$ approximate $E_q^{\star}$.
